\section{Rooted tree gauge and Haar disintegration}\label{app:tree-gauge}

This appendix proves the measure statement used in
\eqref{eq:haar-factorization}.  Let $\Gamma=(V,E)$ be any finite connected
oriented graph, fix a root $o$, and let $T$ be a spanning tree.  The based
gauge group is $\mathcal G_o=\{g\in G^V:g_o=1\}$ for a compact group $G$.

\begin{lemma}[Unique tree representative]\label{lem:unique-tree}
Every $U\in G^E$ has a unique based gauge transformation $g\in\mathcal G_o$
such that $(g\cdot U)_e=1$ on every tree edge.
\end{lemma}

\begin{proof}
Orient the tree away from the root.  Set $g_o=1$.  If $e=(x,y)$ is the tree
edge from a vertex whose $g_x$ has already been defined to a new vertex, the
condition $g_xU_{xy}g_y^{-1}=1$ forces $g_y=g_xU_{xy}$.  Induction over tree
distance defines every $g_y$ and proves uniqueness.
\end{proof}

\begin{proposition}[Exact product-Haar factorization]\label{prop:tree-haar}
Let $\dd U_E$ and $\dd g_{V\setminus\{o\}}$ denote normalized product Haar
measure.  Under the bijection
\[
 G^E\longleftrightarrow
 \mathcal G_o\times G^{E\setminus T},
 \qquad U\longmapsto(g,U^{\mathrm{tree}}|_{E\setminus T}),
\]
product Haar measure factors as
\[
 \dd U_E=\dd g_{V\setminus\{o\}}\,\dd U_{E\setminus T}.
\]
\end{proposition}

\begin{proof}
Order the nonroot vertices so that every parent precedes its children.  For a
leaf $y$ with parent $x$, replace the tree-link variable $U_{xy}$ by
$g_y=g_xU_{xy}$.  Conditional on previously transformed variables, left
translation by $g_x$ preserves Haar measure, so $\dd U_{xy}=\dd g_y$.  The
change sends every cotree edge incident to $y$ to a left or right translate of
its old variable; those translations also preserve Haar measure.  Eliminate
leaves inward until all tree links have been replaced by nonroot gauge
variables.  Fubini's theorem gives the product factorization.  The argument is
finite and introduces no determinant.
\end{proof}

For $G=\SUtwo$, the exponential map in the ball $\norm A<\pi$ pulls normalized
Haar measure back, up to a constant, to
\[
 j(A)\dd A,\qquad j(A)=\left(\frac{\sin\norm A}{\norm A}\right)^2.
\]
Because tree links are fixed to the identity before logarithmic coordinates
are introduced, the reduced chart has exactly the product of these densities
over cotree links.

The quotient is based.  A constant transformation not equal to the identity
at the root is not divided out; it acts on the reduced variables by
simultaneous conjugation.  This avoids stabilizer singularities and is the
boundary convention used throughout the paper.
